\documentclass[stds]{apa7}
\usepackage{amsmath}
\usepackage{amssymb}
\usepackage{graphicx}
\usepackage{booktabs}
\usepackage{biblatex}
\usepackage{csquotes}
\usepackage{hyperref}
\usepackage{url}
\usepackage[utf8]{inputenc}
\usepackage[T1]{fontenc}

\addbibresource{refs.bib}

\graphicspath{{./figures/}}

\title{The Role of Social Media in Academic Performance: A Meta-Analytic Review}

% Note: APA 7 requires a very long abstract
\abstract{
Social media platforms have become ubiquitous in daily life, particularly among
young adults and adolescents. As educational institutions increasingly incorporate
social media into teaching and learning practices, questions arise regarding the impact
of social media use on academic performance. This meta-analysis synthesizes findings
from 127 studies conducted between 2010 and 2024, encompassing data from over 200,000
students across 35 countries. Results indicate a small but significant negative
relationship between social media use and academic performance ($r = -.12$, 95\% CI
$[-\,0.14, -0.10]$). However, this relationship is moderated by several factors
including the type of social media platform, purpose of use (academic versus
recreational), and individual differences in self-regulation. Findings have
important implications for educators and policymakers seeking to promote effective
use of technology in educational contexts.}

\begin{document}

\maketitle

\section{Introduction}

The proliferation of social media has fundamentally changed how students interact,
communicate, and access information.

\section{Method}

We conducted a comprehensive literature search across PsycINFO, Web of Science,
and ERIC databases.

\section{Results}

\begin{table}[ht]
\caption{Summary of meta-analytic findings by platform type}
\label{tab:platforms}
\begin{tabular}{lcccc}
  \toprule
  Platform Type & $k$ & $N$ & $r$ & 95\% CI \\
  \midrule
  Social networking & 45 & 45,000 & -.14 & [-.17, -.11] \\
  Microblogging & 28 & 28,000 & -.09 & [-.12, -.06] \\
  Video platforms & 32 & 32,000 & -.05 & [-.08, -.02] \\
  Academic networks & 22 & 22,000 & .08 & [.04, .12] \\
  \bottomrule
\end{table}
\end{table}

\section{Discussion}

The findings suggest that while social media use is associated with somewhat
lower academic performance, the relationship is nuanced and depends on how
students use these platforms.

\printbibliography

\end{document}
